\documentclass{situationist-preprint}

\title{Wittgenstein's Ladder}
\msid{LNG-22-b}
\msfamily{LNG-22}
\msdiscipline{Philosophy of Language}
\mstype{Consequence}
\mskeywords{Wittgenstein; Tractatus; saying and showing; nonsense; ladder; resolute reading; metalanguage}

\begin{document}
\maketitle

\begin{abstract}
At the end of the \emph{Tractatus}, Wittgenstein says that whoever understands him recognizes its propositions as nonsensical and must throw away the ladder after climbing up it. We derive from the thesis of LNG-22, that some things can be shown but not said though far fewer than the \emph{Tractatus} supposed, what becomes of the ladder. Most of its rungs, those concerning the logical form and semantics of particular languages, are sayable in a richer language and need not be thrown away. One rung cannot be kept as a statement: the generalization over all languages. Any attempt to state it either restricts its scope, and then is a truth that can be kept, or does not restrict it, and then fails. The ladder is shortened, not discarded.
\end{abstract}

\section{The ladder}

\emph{Tractatus} 6.54 says that the book's propositions serve as elucidations: whoever understands their author recognizes them as nonsensical, having used them to climb beyond them, and must throw away the ladder \citep{Wittgenstein1922}. Two readings dominate. On the ineffability reading, the propositions are nonsensical attempts to say what can only be shown, and there are truths, about logical form and the relation of language to the world, that they gesture at \citep{Hacker2000}. On the resolute reading, the propositions are plain nonsense, nothing is shown that they fail to say, and the book's aim is to bring the reader to recognize this \citep{Diamond1988}.

LNG-22-a defends a thesis between these. Some things are shown and not said: the general semantic fact relating every language to what its sentences say (Theorem 3.2 of LNG-22-a). But the logical form and semantics of any particular language are sayable (Lemma 3.1), and so the class of what is shown but not said is much smaller than the \emph{Tractatus} supposed (Corollary 3.3). We derive what this implies for the ladder.

\section{Which rungs remain}

Call a \emph{rung} any proposition of the \emph{Tractatus} about language, logic or the relation between language and world, taken at the generality that its words suggest. The rungs fall into two kinds.

\begin{proposition}\label{p:particular}
Every rung whose subject is the logical form or semantics of a particular language, or of every language in a class not including the language in which it is stated, is said by some sentence of some language.
\end{proposition}

\begin{proof}
By Lemma 3.1 of LNG-22-a, the semantics and logical form of each particular language are sayable in a richer metalanguage. A claim about every language in a class that excludes the metalanguage is a claim about languages whose semantics that metalanguage can state, and it is stated there.
\end{proof}

The suggestion is \citeauthor{Russell1922}'s, in his introduction to the \emph{Tractatus}: that what cannot be said about the structure of one language can be said in another, and so on up a hierarchy \citep{Russell1922}. Proposition~\ref{p:particular} vindicates it for every rung that concerns any particular language, including the picture theory applied to a given language, the account of truth functions, and the claim that logical propositions are tautologies. These rungs are not nonsense. They need not be thrown away; they need only be relocated to a language that is not the one they describe.

\begin{proposition}\label{p:top}
A rung that states the general semantic fact, for every language including its own, is either restricted in scope, and then is a truth sayable by Proposition~\ref{p:particular}, or unrestricted, and then is said by no sentence.
\end{proposition}

\begin{proof}
Suppose a sentence of a language $M$ purports to state the semantic relation of every language to what its sentences say. If its quantifier excludes $M$, Proposition~\ref{p:particular} applies. If it includes $M$, then $M$ can state semantic facts about every language, hence describes its own syntax by Assumption 2.4 of LNG-22-a, and contains a truth predicate for its own sentences, contrary to Assumption 2.3. So the sentence does not state what it purports to state.
\end{proof}

The top rung is therefore the only one to be thrown away, and what is thrown away is not a truth but a sentence that fails to express the truth it aims at. What it aims at is shown, in the sense of Definition 2.2 of LNG-22-a: competence in each language consists partly in using it in accordance with the general semantic fact.

\begin{objection}
On the resolute reading, the reader who has climbed the ladder is left with nothing but the recognition that she was taken in. The intermediate position keeps a shown fact, and so is the ineffability reading under another name.
\end{objection}

\begin{reply}
It differs from the ineffability reading in what it places beyond saying, and from the resolute reading in what it calls nonsense. The ineffability reading puts the logical form of every proposition, the limits of the world and ethics beyond saying; Corollary 3.3 of LNG-22-a removes all of these. The resolute reading calls every rung nonsense; Proposition~\ref{p:particular} shows that most rungs are true statements about particular languages. What remains of 6.54 is one claim: that the generalization over all languages, read without restriction, is not a statement. LNG-22-a answers the related objection that what cannot be said cannot be whistled either.
\end{reply}

\section{Conclusion}

The ladder of the \emph{Tractatus} is mostly sound. Its rungs about particular languages can be kept, stated from a richer standpoint; only the rung that generalizes over every language, its own included, must be given up as a statement, and what it was meant to state is shown in the use of language. Wittgenstein was right that something must be thrown away and wrong about how much.

\begin{thebibliography}{9}\small
\bibitem[Diamond(1988)]{Diamond1988} Diamond, C. (1988). Throwing away the ladder. \emph{Philosophy}, 63(243), 5--27.
\bibitem[Hacker(2000)]{Hacker2000} Hacker, P.~M.~S. (2000). Was he trying to whistle it? In A. Crary and R. Read (eds), \emph{The New Wittgenstein}. Routledge.
\bibitem[Russell(1922)]{Russell1922} Russell, B. (1922). Introduction. In L. Wittgenstein, \emph{Tractatus Logico-Philosophicus}. Kegan Paul.
\bibitem[Wittgenstein(1922)]{Wittgenstein1922} Wittgenstein, L. (1922). \emph{Tractatus Logico-Philosophicus}. Trans. C.~K. Ogden. Kegan Paul. Cited by proposition number.
\end{thebibliography}

\end{document}
